\section{Height representation and critical inputs}\label{sec:heights}

We first pass from spins to Bessel heights, fixing the normalization
and the boundary conventions.  We then record the finite comparisons
and critical height limits used in the geometric and analytic arguments.
These inputs concern the Bessel weights obtained from the XY model.

\subsection{Dual partition functions}

Choose one orientation of each edge of a finite nearest-neighbor
graph \(G=(V,E)\) and put \(\nabla_eh=h_y-h_x\) for \(e=(x,y)\).
We call these height vertices and edges \emph{primary}, to distinguish
them from the auxiliary subdivisions in the comparison proof;
\emph{primal} refers to the spin lattice.
The height takes values in \(2\pi\Z\).  Its increment weight is
\begin{equation}\label{eq:bessel-weight}
 p_b(j)=e^{-b}I_j(b),\qquad
 \sum_{j\in\Z}p_b(j)e^{tj}=\exp\{b(\cosh t-1)\}.
\end{equation}
Here \(I_j\) is the modified Bessel function.  Thus \(p_b\) is the law
of the difference of two independent Poisson variables of mean \(b/2\).
An \emph{integral connection} is an antisymmetric edge function
\(\tau:E\to2\pi\Z\).  The weight in that sector is
\[
 \prod_{e\in E}p_b\left(\frac{\nabla_eh+\tau_e}{2\pi}\right).
\]
A connection is flat where its circulation around every contractible
cycle is zero.
Its circulations around the remaining holes are its periods.

We allow nonnegative quadratic penalties in real linear observations,
hard pins, and prescribed relative height data.  Every observation
used after a cut is supported in one component.  A component with no
absolute constraint or confining observation is summed modulo a common
translation by \(2\pi\Z\); when its constant is confined, all its
translation copies are summed.  A hard pin fixes that constant instead.
Relative translations between cut components are summed again on
restoration.  Write \(Z_G^\tau\) for a partition function with these
conventions and \(Z_G^0\) for the same graph, cuts, and penalties with
all affine data centered.

An integral gauge change replaces \(\tau\) by \(\tau+\nabla\psi\) and
the height variable by \(h-\psi\), where \(\psi:V\to2\pi\Z\).
This is an exact reindexing when every affine datum is changed with
the height.  A penalty \(\|Bh-y\|^2/2\) then has center \(y-B\psi\),
and absolute data \(b^p\) on a set \(p\) become \(b^p-\psi|_p\).
For relative data \(h_z-h_{z_0}=u_z\), the new value is
\(u_z-\psi_z+\psi_{z_0}\).  Thus a hard zero pin is preserved when
\(\psi=0\) on it, and a constraint requiring all heights on a set
to be equal is preserved when \(\psi\) is constant there.

\begin{lemma}[Spin--height duality]\label{lem:duality}
In a free spin box, the two-point correlation is \(Z_G^\tau/Z_G^0\)
on the dual graph with the exterior face fixed at height zero.
One may represent that face by separate exterior incident copies, all
pinned at zero, or identify those copies while retaining every edge.
Circulations below are read in the dual with the exterior copies identified.
The connection has circulations \(+2\pi\) and \(-2\pi\) around the
two sources and can be supported on the dual bonds crossing a primal
current path joining them.

The center magnetization \(M_n\) defined in
\eqref{eq:center-definition} satisfies
\(M_n=\E_{n,b_c}^{\,0}e^{i\theta_0}>0\).
It is \(Z_G^\tau/Z_G^0\) on the free dual square with
\(2n\) dual vertices on each side, indexed by the primal cells, where
\(\tau\) has one circulation \(2\pi\) around the center.  It may be
supported on the dual bonds crossing a primal current path from the
center to the boundary.
\end{lemma}

\begin{proof}
Expand each spin interaction in its Fourier series,
\[
 e^{b\cos(\theta_x-\theta_y)}
 =\sum_{j\in\Z}I_j(b)e^{ij(\theta_x-\theta_y)}.
\]
In the fixed-boundary square first cancel the interactions whose two
endpoints are fixed boundary spins.  Integrating each remaining
unfixed spin imposes its integer divergence condition.  Subtract a
fixed primal current carrying the inserted charges along the stated
path.  The remaining current is divergence free at every integrated
vertex: at all vertices in the free spin box and at the interior
vertices in the fixed-boundary square.  Rotation onto dual edges then
gives an integral gradient.  In the free spin box its exterior height
is fixed; in the fixed-boundary square it is determined modulo a common
translation.  The factors \(e^b\) in \eqref{eq:bessel-weight} cancel
in the ratio.  The numerator and denominator are sums of strictly
positive Bessel weights, so the ratios are positive.  Current reversal
identifies each insertion with its conjugate, giving the stated real
cosine expectations.  This is the duality of \cite[Lemma~2.1]{C} and
\cite[Lemma~2.1]{S}, with their respective boundary conventions.
For the classical current and Bessel-height representations, see also
\cite[Appendix~A]{FrohlichSpencer1981} and
\cite[Section~2]{vanEngelenburgLis2023}.
\end{proof}

The reference form on a retained primary graph is the unit-conductance
form
\begin{equation}\label{eq:reference-energy}
 \ip{g}{A_Gg}=\sum_{\{x,y\}\in E}(g_x-g_y)^2.
\end{equation}
It converges under lattice scaling to \(\int|\nabla g|^2\).
For the four nearest steps this is exactly the operator \(A_0\) in
\cite{B}.  We write \(A\) when the graph is understood.

\subsection{The critical height inputs}

\begin{proposition}[Finite affine comparisons]\label{prop:comparisons}
On every fixed primary graph, the Bessel height law has the following
properties.
\begin{enumerate}
\item A shifted-to-centered partition ratio is at most one and
decreases when the same nonnegative quadratic precision is added to
the affine and centered coordinates.  Affine data may include an
integral connection, observation centers, and compatible nonzero
equality values.  The monotonicity remains valid with a fixed real
linear statistic in the numerator exponent, although the tilted
ratio need not be at most one.
\item For zero affine data, centered covariances and centered real moment generating
functions decrease under added centered precision.  Centered
characteristic coefficients allowed by the constant convention in both
laws are nonnegative and increase under added precision.  They are
strictly positive for the finite-graph phases
used below, whose coefficient total on every unconfined component is
an integer.
\item If \(X=\sum_e q_e\nabla_eh\) in a centered law, then
\[
 \Var X\le C_b\sum_eq_e^2,\qquad
 \E|X|^k\le C_{k,b}\left(\sum_eq_e^2\right)^{k/2}.
\]
More generally, after releasing endpoint, curl, and period constraints
while retaining independent primary increments, a tilt with
\(\max_e|q_e|\le K\) costs at most
\(\exp(C_{b,K}\sum_eq_e^2)\).
\end{enumerate}
The statements include independent Gaussian observation noises.
Released constants are taken modulo translations or confined; a real
statistic used in releasing a constant must annihilate it.
\end{proposition}

\begin{proof}
The comparison assertions are \cite[Lemma~2.2]{C}.  That lemma replaces
each primary bond by a chain of lattice-Gaussian bonds, applies the
finite affine comparison on the enlarged lattice, and passes to the
limit at the fixed primary graph.  Equality constraints are limits
of positive quadratic penalties.  Thus the same affine statistic is
kept when observations or constraints are released.
Strict positivity for the phases used here follows from the
phase-transfer estimate \cite[Equation~(2.2) and its proof]{C}: on an
unconfined component subtract an integer point charge and interpolate
the remaining zero-total finite-variance test from zero; on a confined
component interpolate directly using its finite fixed-graph variance.
The finite lattice-Gaussian comparison underlying this argument is
developed in \cite[Lemma~2.1]{B} from the inequality of
\cite[Theorem~2.1 in the arXiv version]{RSD2017}.

For the last assertion, perform these releases before estimating the
tilt.  The independent primary transform is the product of the
transforms in \eqref{eq:bessel-weight}, with \(t=2\pi q_e\).
For \(|q_e|\le K\),
\(\cosh(2\pi q_e)-1\le C_Kq_e^2\).  This gives the last bound.
Normalize \(q\) by \((\sum q_e^2)^{1/2}\) and differentiate bounded
exponential moments to obtain the stated moment bounds.  These are
the primary reference-cost bounds of \cite[Lemma~2.3]{C}; no bound
on the microscopic Gaussian precision of the subdivided chains is
used.
\end{proof}

The next proposition uses absolute pin data: \(h_H=0\) fixes every
primary height in \(H\), and \(b^p\) specifies a compatible absolute
height at every primary vertex of \(p\).  Its approximation by finite
guard observations is understood in probability, at each prescribed
tolerance, under the sampled law stated there.  The finite list of
smooth guard probes is fixed before the lattice limit and enlarged
afterward to reduce the tolerance.

\begin{proposition}[Critical coefficient and mixed limits]\label{prop:height-inputs}
For the critical parameter in \eqref{eq:critical-definition}, the
Bessel height coefficient of \cite{C} is
\begin{equation}\label{eq:critical-coefficient}
                         a=a(b_c)=8\pi.
\end{equation}
At this parameter the following fixed-geometry statements hold.
\begin{enumerate}
\item Let the fixed geometry belong to the polygonal Lipschitz class
of \cite[Theorem~3.1]{C}: it is a finite union of bounded rectangular
pieces with nondegenerate faces and widths.  Pieces in one component
are joined along full interfaces, which may be subdivided into segments;
point contacts do not join components.  All passages have positive width.
Allow free walls and
hard pins on the primary vertices in aligned unions of positive-area
boxes.  On each component
\(U\), fixed bounded, piecewise continuous density smears of the height
have Gaussian Laplace and moment limits
with covariance the inverse of
\[
 a^{-1}\int_U|\nabla v|^2
 \quad\text{on}\quad
 \{v\in H^1(U):v=0\text{ on the pin}\}.
\]
For an unpinned component this is a statement about mean-free smears;
its fractional mean modulo \(2\pi\) becomes independent and uniform.
Fixed finite-rank nonnegative penalties in box averages give the
corresponding Gaussian laws, with constants integrated when the
penalties confine them.  The statements hold along every admissible
mesh sequence, including bounded rounding of the fixed geometry.
They include energy-approximable shifts in the sense of
\cite[Theorem~3.1]{C}.
\item Start with no hard pins and use the centered law whose
component-local box-average penalties confine every pre-pin component
constant.  Let \(H,p\) be positively separated unions of boxes.  Put
\[
 D(H,p)=\log\frac{\Prob(h_H=0\mid h_p=0)}{\Prob(h_H=0)},\qquad
 r_H(b^p)=\frac{\Prob(h_H=0\mid h_p=b^p)}
                 {\Prob(h_H=0\mid h_p=0)}.
\]
Then \(D(H,p)\ge0\) converges to the Gaussian pin interaction,
\(0<r_H\le1\), and
\(\E_{\mathrm{cen}} r_H=e^{-D(H,p)}\), where the expectation samples
the guard data in the centered law.
For every fixed finite list of these density smears and box averages,
the conditional Laplace transforms and \(r_H\) admit approximations
by fixed finite lists of smooth observations of the guard data, whose
laws and functions converge to their Gaussian counterparts.  The
list is enlarged after the lattice limit.  These convergence assertions
also hold when the data are sampled under a fixed finite real linear
tilt; \(D\) and \(r_H\) retain their centered definitions above.
\end{enumerate}
\end{proposition}

\begin{proof}
Equation~\eqref{eq:critical-coefficient} is the conclusion of
\cite[Section~6]{C}.  The first and second parts are respectively
\cite[Theorems~3.1 and~3.2]{C}.  In the second part, observations are
component-local before pins are imposed.  The finite-observation
formulation is part of the theorem; it does not assert uniform
convergence at every microscopic guard value.
\end{proof}

We also use the restoration consequence \cite[Lemma~3.9]{C}.
Suppose every endpoint of a restored bond, and the support of every
restored form, lies in a core \(H\) separated from a guard \(p\).
If \(r_\Delta\) is the conditional restoration factor divided by
its value at zero guard data, then \(r_H\le r_\Delta\le1\).
The guard laws in the two topologies transfer events whose
probabilities tend to zero in either direction, also under a fixed
real linear tilt.  The finite-observation approximations above concern
conditional Laplace transforms and \(r_H\); no additional limit for
\(r_\Delta\) is asserted.  In each use, the graph, pin regions, and
observation cells are fixed before the lattice limit.
